\documentclass[10pt,a4paper]{amsart}
\usepackage[margin=19mm]{geometry}
\usepackage{amsmath,amssymb,mathtools}
\usepackage[scr=rsfs,cal=euler]{mathalfa}
\usepackage{xfrac}
\usepackage[unicode,hidelinks,bookmarks=false]{hyperref}
\newcommand{\ie}{i.e.\ }
\newcommand{\cf}{cf.\ }
\newcommand{\resp}{respectively\ }
\newcommand{\Subsection}{\subsection}
\providecommand{\nobreakdash}{\nobreak\hbox{-}}
\setlength{\emergencystretch}{2em}
\multlinegap0pt
\usepackage{bm}
\usepackage{bbm}
\usepackage{dsfont}
\usepackage{xspace}
\usepackage{cancel}
\usepackage{mathtools}
\usepackage{amsthm}
\makeatletter
\@ifundefined{theorem}{\newtheorem{theorem}{Theorem}}{}
\@ifundefined{definition}{\newtheorem{definition}[theorem]{Definition}}{}
\@ifundefined{lemma}{\newtheorem{lemma}[theorem]{Lemma}}{}
\makeatother
\makeatletter
\expandafter\ifx\csname manualtheorem\endcsname\relax
\newenvironment{manualtheorem}[1]{%
\renewcommand\themanualtheoreminner{#1}%
  \manualtheoreminner
}{\endmanualtheoreminner}
\fi
\makeatother
\title[Global dimension of a string algebra]{Global dimension of a string algebra}
\author{Zheng Xin, Lingchun Zhang}
\date{Source: arXiv:2607.11552v2 (CC BY 4.0)}
\begin{document}
\maketitle

\noindent\emph{Source and licence.} Global dimension of a string algebra. Version arXiv:2607.11552v2 (CC BY 4.0), distributed under the Creative Commons Attribution 4.0 International licence, as recorded in the arXiv licence metadata for this exact version and shown on the abstract page https://arxiv.org/abs/2607.11552v2. Full rights notice: licenses/arxiv-2607.11552-CC-BY-4.0.txt.\par\smallskip

\noindent\textit{Editorial scope.} Counted unit: the Lemma whose source label is \texttt{lem:proj-formula} in the article (source0.tex lines 424--426, source label \texttt{lem:proj-formula}), delivered together with the article's own complete original proof of it (source0.tex lines 427--583, one \texttt{proof} environment of 157 lines). No mathematics is written, repaired or re-derived: statement and proof are the article's own text line for line. Required context retained uncounted: equ:proj (lines 378--412). Every definition, display and label of those lines is retained unchanged. Omitted from this package and not redistributed: 1--377, 413--423, 584--946. Nothing inside a retained excerpt is omitted or altered: every retained body is delivered exactly as the article's lines are written, so no omission receipt of any kind accompanies this package. Citations: the retained excerpts carry no citation command at all, so no bibliography entry of the article is transcribed and no reference is redistributed. Reference adaptation: none was needed and none was made; every reference inside the delivered unit resolves inside it, so no unresolved reference or citation is produced. Printed slips kept literal: no printed word, symbol, hypothesis or number of the counted statement or of its proof was altered. Layout and class: the article's own class is \texttt{amsart} with packages including amssymb, latexsym, amsmath, amscd, amsthm, amsfonts, enumerate, multirow, color, xy, etex, caption, soul, float; the wrapper is the lane's pinned \texttt{amsart} editorial wrapper at 10pt on A4, which restates the theorem-like environments the retained text needs and adds the article's own macro definitions that the retained text needs (none), with extra wrapper packages (bm, bbm, dsfont, xspace, cancel, mathtools). The delivered numbering is the unit's own. Assets: no retained span contains a figure environment, a graphics-inclusion command, a PDF-page inclusion, a table or a TikZ picture; no image file is copied into this package, the package declares no asset dependency and no image file is redistributed with it. Licence: version arXiv:2607.11552v2 of this article is distributed under the Creative Commons Attribution 4.0 International licence, as recorded in the arXiv licence metadata for this exact version and shown on the abstract page https://arxiv.org/abs/2607.11552v2. A full rights notice is delivered at licenses/arxiv-2607.11552-CC-BY-4.0.txt. Characters: every retained excerpt is pure ASCII, so the delivered engine, XeLaTeX with its default Latin Modern text font, reports no missing character.
\begin{definition}
	For any vertex \( v \in Q_0 \), define a resolution \( P^\bullet \to S(v) \):
	
	\begin{equation} \label{equ:proj}
		\cdots \xrightarrow{\partial^{-i - 1}} P^{-i} \xrightarrow{\partial^{-i}} P^{-i + 1} \xrightarrow{\partial^{-i + 1}} \cdots \xrightarrow{\partial^{-2}} P^{-1} \xrightarrow{\partial^{-1}} P^0 \xrightarrow{\partial^0} S(v) \to 0  \tag{$\divideontimes$}.
	\end{equation}
	

	The modules in the resolution are:
\begin{gather*}
	P^0 = P(v), \quad P^{-1} = \bigoplus_{\alpha \in s^{-1}(v)} P(t(\alpha)), \\
	P^{-i} = \bigoplus_{\substack{\alpha \in s^{-1}(v) \\ W \in \mathsf{Rel}(\alpha)_i}} P(t(\mathscr{L}(W))), \quad i > 1.
\end{gather*}
	where \( W = w_1 \cdot w_2 \cdots w_n \), if \( \alpha \) does not exist, then \( P(t(\alpha)) = 0 \); If \( W \) does not exist, then \( P(t(\mathscr{L}(W))) = 0 \).
	
	The differentials in the resolution are:
	
	$\partial^0 = \pi$,
	 
	$\partial^{-1} = (p(\delta)~~p(\epsilon))$, $\delta \neq \epsilon$, $s(\delta) = s(\epsilon) = v$, if  $\delta$  does not exist, then  $p(\delta) = 0$, similarly for  $\epsilon$,
	
	For \( i > 1 \), \( \partial^{-i} = (\partial^{-i}_{jk})_{1 \leq j \leq 2^{i - 1}, 1 \leq k \leq 2^i} \), where
	\[
	\partial^{-i}_{jk} =
	\begin{cases}
		p(\mathscr{L}(W)), \, j = \mathcal{L}(p(w_{i - 1})), \, k = 2j - 1, \, \mathscr{L}(W) \text{ is the short kernel of } p(w_{i - 1}) \\
		\substack{\alpha \in s^{-1}(v) \\ W \in \mathsf{Rel}(\alpha)_i} \\ \\
		p(\mathscr{L}(W)), \, j = \mathcal{L}(p(w_{i - 1})), \, k = 2j, \, \mathscr{L}(W) \text{ is the long kernel of } p(w_{i - 1}) \\
		\substack{\alpha \in s^{-1}(v) \\ W \in \mathsf{Rel}(\alpha)_i} \\  \\
		0, \, \text{otherwise}
	\end{cases}
	\]
	
	Here, \( \pi \) is the natural projection, \( S(v) \) and \( P(v) \) are the simple module and projective module at the vertex \( v \) respectively. The function $\mathcal{L}$ is used to calculate the column index of elements in a matrix.
\end{definition}

\begin{lemma}\label{lem:proj-formula}
	Let \( A = \mathbf{k}Q/I \) be a string algebra, \( v \in Q_0 \), then the resolution (\ref{equ:proj}) is the minimal projective resolution of the simple module \( S(v) \).
\end{lemma}

\begin{proof}
	Since 
	\begin{center}
		\( P^0 = P(v) \), \( P^{-1} = \bigoplus\limits_{\alpha \in s^{-1}(v)} P(t(\alpha)) \), 
		
		\( P^{-i} = \bigoplus\limits_{\substack{\alpha \in s^{-1}(v) \\ W \in \mathsf{Rel}(\alpha)_i}} P(t(\mathscr{L}(W))) \) (\( i > 1 \)),
	\end{center}
	 for any \( i \geq 0 \), \( P^{-i} \) is a projective module.
	
	For \( i > 1 \), from the definition of \( \partial^{-i} = (\partial^{-i}_{jk})_{1 \leq j \leq 2^{i- 1 }, 1 \leq k \leq 2^i} \), we have:
	\[
	\partial^{-i}_{jk} = \begin{pmatrix}
		p(u_1) & p(u_2) & 0 & 0 & 0 & 0 & \cdots & 0 & 0 \\
		0 & 0 & p(u_3) & p(u_4) & 0 & 0 & \cdots & 0 & 0 \\
		\vdots & \vdots & \vdots & \vdots & \vdots & \vdots & \ddots & \vdots & \vdots \\
		0 & 0 & 0 & 0 & 0 & 0 & \cdots & p(u_{2^i - 1}) & p(u_{2^i})
	\end{pmatrix}_{2^{i- 1 } \times 2^i}
	\]
	In the \( j \)-th (\( 1 \leq j \leq 2^{i- 1 } \)) row of the differential \( \partial^{-i}_{jk} \), the possible non-zero elements are \( p(u_{2j - 1}) \) and \( p(u_{2j}) \), which are located in the \( (2j - 1) \)-th column and the \( 2j \)-th column respectively. Here, \( u_1, u_2, \ldots, u_{2^i - 1}, u_{2^i} \) correspond one-to-one to the elements \( W_1, W_2, \ldots, W_{2^i - 1}, W_{2^i} \) in \( \mathsf{Rel}(\alpha)_i \), and \( \mathscr{L}(W_m) = u_m \), \( 1 \leq m \leq 2^i \), \( m \in \mathbb{Z} \).
	\[
	\partial^{-i}: P^{-i} = \bigoplus\limits_{\substack{\alpha \in s^{-1}(v) \\ W \in \mathsf{Rel}(\alpha)_i}} P(t(\mathscr{L}(W))) \longrightarrow P^{-i + 1} = \bigoplus\limits_{\substack{\alpha \in s^{-1}(v) \\ W \in \mathsf{Rel}(\alpha)_{i - 1}}} P(t(\mathscr{L}(W))).
	\]
	For \( (m_1, m_2, \ldots, m_{2^i - 1}, m_{2^i})^T \in P^{-i} \), where \( m_1, m_2, \ldots, m_{2^i - 1}, m_{2^i} \) belong to different direct summands of \( P^{-i} \), we have:
	
	\[
	\begin{aligned}
		&\partial^{-i} \cdot (m_1, m_2, \ldots, m_{2^i - 1}, m_{2^i})^T \\
		&= \begin{pmatrix}
			p(u_1) & p(u_2) & 0 & 0 & 0 & 0 & \cdots & 0 & 0 \\
			0 & 0 & p(u_3) & p(u_4) & 0 & 0 & \cdots & 0 & 0 \\
			\vdots & \vdots & \vdots & \vdots & \vdots & \vdots & \ddots & \vdots & \vdots \\
			0 & 0 & 0 & 0 & 0 & 0 & \cdots & p(u_{2^i - 1}) & p(u_{2^i})
		\end{pmatrix} \begin{pmatrix}
			m_1 \\
			m_2 \\
			\vdots \\
			m_{2^i}
		\end{pmatrix} \\
		&= \begin{pmatrix}
			p(u_1)(m_1) + p(u_2)(m_2) \\
			p(u_3)(m_3) + p(u_4)(m_4) \\
			\vdots \\
			p(u_{2^i - 1})(m_{2^i - 1}) + p(u_{2^i})(m_{2^i})
		\end{pmatrix} \\
		&= \begin{pmatrix}
			u_1 m_1 + u_2 m_2 \\
			u_3 m_3 + u_4 m_4 \\
			\vdots \\
			u_{2^i - 1} m_{2^i - 1} + u_{2^i} m_{2^i}
		\end{pmatrix}.
	\end{aligned}
	\]
	
	Therefore, 
	\begin{center}
	\begin{align*}
		\text{Im}(\partial^{-i}) 
		&= (u_1 A \oplus u_2 A) \oplus (u_3 A \oplus u_4 A) \oplus \cdots \oplus (u_{2^i - 1} A \oplus u_{2^i} A) \\
		&= \bigoplus\limits_{\substack{\alpha \in s^{-1}(v) \\ W \in \mathsf{Rel}(\alpha)_i}} \mathscr{L}(W)A .
	\end{align*}
	\end{center}
	
	If \( (m_1, m_2, \cdots, m_{2^i - 1}, m_{2^i})^T \in \ker (\partial^{-i}) \), then we have
	\[
	\begin{cases}
		u_1 m_1 + u_2 m_2 = 0 \\
		u_3 m_3 + u_4 m_4 = 0 \\
		\vdots \\
		u_{2^i - 1} m_{2^i - 1} + u_{2^i} m_{2^i} = 0
	\end{cases}.
	\]
	Since \( u_1 m_1 \) and \( u_2 m_2 \) are linearly independent, we know \( u_1 m_1 = 0 \) and \( u_2 m_2 = 0 \). Similarly, we can obtain:
	\[
	u_3 m_3 = u_4 m_4 = \cdots = u_{2^i - 1} m_{2^i - 1} = u_{2^i} m_{2^i} = 0.
	\]
	That is, 
	\begin{center}
		\( m_1 \in \ker p(u_1) \), \( m_2 \in \ker p(u_2) \), \(\ldots\), \( m_{2^i - 1} \in \ker p(u_{2^i - 1}) \), \( m_{2^i} \in \ker p(u_{2^i}) \).
	\end{center} 
	Thus,
   \[
	\begin{aligned}
		\ker(\partial^{-i})  &\subseteq \ker p(u_1) \oplus \ker p(u_2) \oplus \cdots \oplus \ker p(u_{2^i - 1}) \oplus \ker p(u_{2^i}) \\
		&= \bigoplus\limits_{\substack{\alpha \in s^{-1}(v) \\ W \in \mathsf{Rel}(\alpha)_i}} \ker p(\mathscr{L}(W))\\ 
		& = \bigoplus\limits_{\substack{\alpha \in s^{-1}(v) \\ W \in \mathsf{Rel}(\alpha)_{i + 1}}} \mathscr{L}(W)A
	\end{aligned}.
	\]
	
	It is clear that 
\begin{center}
	\( \ker p(u_1) \oplus \ker p(u_2) \oplus \cdots \oplus \ker p(u_{2^i - 1}) \oplus \ker p(u_{2^i}) \subseteq \ker(\partial^{-i})  \). 
\end{center}
Therefore, 
     \[
		\ker(\partial^{-i}) = \bigoplus\limits_{\substack{\alpha \in s^{-1}(v) \\ W \in \mathsf{Rel}(\alpha)_{i + 1}}} \mathscr{L}(W)A. 
\] 
	Hence, for \( i > 1 \), \( \text{Im}(\partial^{-(i + 1)}) = \ker(\partial^{-i}) \).
	
	Next, we prove 
	\( \text{Im}(\partial^{-2}) = \ker(\partial^{-1}) \) and \( \text{Im}(\partial^{-1}) = \ker(\pi) \).
	
	It is known that \( \text{Im}(\partial^{-1}) = \bigoplus\limits_{\alpha \in s^{-1}(v)} \alpha A \), which is the submodule generated by all paths greater than 0 of \( P^0 = P(v) \), that is, \( \text{Im}(\partial^{-1}) = \ker(\pi) \). In addition, if \( v \) is the starting point of two different arrows \( \delta, \epsilon \), then \( \partial^{-1} = (p(\delta)~~~~p(\epsilon)) \). 
	For any element \( (x, y)^T \in \ker(\partial^{-1}) \), where \( x \in P(t(\delta)) \), \( y \in P(t(\epsilon)) \), then \( \partial^{-1}(x, y)^T = \delta x + \epsilon y = 0 \). And since \( \delta x, \epsilon y \) are linearly independent, \( \delta x = 0 \), \( \epsilon y = 0 \), thus \( x \in \ker p(\delta) \), \( y \in \ker p(\epsilon) \), that is
	\[
	\begin{aligned}
		\ker(\partial^{-1})  &\subseteq \ker p(\delta) \oplus \ker p(\epsilon) \\
		&= \bigoplus\limits_{\alpha \in s^{-1}(v)} \ker p(\alpha) \\
		&= \bigoplus\limits_{\substack{\alpha \in s^{-1}(v) \\ W \in \mathsf{Rel}(\alpha)_1}} \ker p(\mathscr{L}(W)) \\
		&= \bigoplus\limits_{\substack{\alpha \in s^{-1}(v) \\ W \in \mathsf{Rel}(\alpha)_2}} \mathscr{L}(W)A \\
		&= \text{Im}~(\partial^{-2}).
	\end{aligned}
	\]
	It is clear that \( \ker p(\delta) \oplus \ker p(\epsilon) \subseteq \ker(\partial^{-1}) \).
    Therefore, \( \ker(\partial^{-1}) = \text{Im}(\partial^{-2}) \).
	
	Next, we need to prove the minimality of resolution \ref{equ:proj}, that is, we need to prove that each differential \( \partial^{-i} (i \geq 1) \) and the natural projection \( \pi \) in the resolution are projective covers. 
	In other words, we need to prove 
    \( \ker(\pi) \subseteq P^0 \text{rad}(A) \)~~and~~~ \( \ker(\partial^{-i}) \subseteq P^{-i} \text{rad}(A) , i \geq 1 \) . 

	For any vertex \( v \in Q_0 \), we have 
	\begin{align*}
		P(v)\operatorname{rad}(A) &= \operatorname{rad}(P(v)) \\
		&= \bigoplus_{\alpha \in s^{-1}(v)} \alpha A .
	\end{align*}
	
	Then, when \( i \geq 2 \),
	\begin{align*}
		P^{-i}\operatorname{rad}(A) 
		&= \bigoplus_{\substack{\alpha \in s^{-1}(v) \\ W \in \mathsf{Rel}(\alpha)_i}} P(t(\mathscr{L}(W)))\operatorname{rad}(A) \\
		&= \bigoplus_{\substack{\alpha \in s^{-1}(v) \\ W \in \mathsf{Rel}(\alpha)_i \\ \beta \in s^{-1}(t(\mathscr{L}(W)))}} \beta A, \\[4pt]
		\ker(\partial^{-i}) 
		&= \bigoplus_{\substack{\alpha \in s^{-1}(v) \\ W \in \mathsf{Rel}(\alpha)_{i + 1}}} \mathscr{L}(W)A \\
		&\subseteq P^{-i}\operatorname{rad}(A).
	\end{align*}
	
	In addition,
	\begin{align*}
		P^{-1}\operatorname{rad}(A) 
		&= \bigoplus_{\alpha \in s^{-1}(v)} P(t(\alpha))\operatorname{rad}(A) \\
		&= \bigoplus_{\substack{\alpha \in s^{-1}(v) \\ \beta \in s^{-1}(t(\alpha)) }} \beta A, \\[4pt]
		\ker(\partial^{-1})
		&= \bigoplus_{\substack{\alpha \in s^{-1}(v) \\ W \in \mathsf{Rel}(\alpha)_2}} \mathscr{L}(W)A \\
		&\subseteq P^{-1}\operatorname{rad}(A).
	\end{align*}
	
	Moreover,
	\begin{align*}
		P^0\operatorname{rad}(A) 
		&= P(v)\operatorname{rad}(A) \\
		&= \bigoplus_{\alpha \in s^{-1}(v)} \alpha A \\
		&= \ker(\pi).
	\end{align*}
	
	Therefore, the resolution \ref{equ:proj} is the minimal projective resolution of the simple module \( S(v) \).
	
	
\end{proof}

\end{document}
